\section{Yang--Mills gradient flow}
\label{sec:3}
\subsection{Flow equations and the perturbative expansion}
The Yang--Mills gradient flow is a deformation of a gauge field configuration
generated by a gradient flow in which the Yang--Mills action integral is
regarded as a potential height. To be explicit, for the gauge
potential~$A_\mu(x)$, the flow is defined
by~\cite{Luscher:2010iy,Luscher:2011bx}
\begin{equation}
   \partial_tB_\mu(t,x)=D_\nu G_{\nu\mu}(t,x)
   +\alpha_0D_\mu\partial_\nu B_\nu(t,x),\qquad
   B_\mu(t=0,x)=A_\mu(x),
\label{eq:(3.1)}
\end{equation}
where $t$~is the flow time and $G_{\mu\nu}(t,x)$ is the field strength of the
flowed field,
\begin{equation}
   G_{\mu\nu}(t,x)
   =\partial_\mu B_\nu(t,x)-\partial_\nu B_\mu(t,x)
   +[B_\mu(t,x),B_\nu(t,x)],
\label{eq:(3.2)}
\end{equation}
and the covariant derivative on the gauge field is
\begin{equation}
   D_\mu=\partial_\mu+[B_\mu,\cdot].
\label{eq:(3.3)}
\end{equation}
The first term in the right-hand side of~Eq.~\eqref{eq:(3.1)} is the
``gradient'' in the functional space, $-g_0^2\delta S/\delta B_\mu(t,x)$, where
$S$~is the Yang--Mills action integral for the flowed field. Note that since
$D_\nu G_{\nu\mu}(t,x)=\Delta B_\mu(t,x)+O(B^2)$, Eq.~\eqref{eq:(3.1)} is a sort
of diffusion equation and the flow for~$t>0$ effectively suppresses
high-frequency modes in the configuration. The second term in the right-hand
side of~Eq.~\eqref{eq:(3.1)} with the parameter~$\alpha_0$ is a ``gauge-fixing
term'' that makes the perturbative expansion well defined. It can be shown,
however, that any gauge-invariant quantities are independent of~$\alpha_0$. In
actual perturbative calculation in the next section, we adopt the ``Feynman
gauge'' $\alpha_0=1$ with which the expressions become simplest.

The formal solution of~Eq.~\eqref{eq:(3.1)} is given
by~\cite{Luscher:2010iy,Luscher:2011bx}
\begin{equation}
   B_\mu(t,x)
   =\int\mathrm{d}^Dy\left[
   K_t(x-y)_{\mu\nu}A_\nu(y)
   +\int_0^t\mathrm{d}s\,K_{t-s}(x-y)_{\mu\nu}R_\nu(s,y)
   \right],
\label{eq:(3.4)}
\end{equation}
where\footnote{Throughout the present paper, we use the abbreviation,
\begin{equation}
   \int_p\equiv\int\frac{\mathrm{d}^Dp}{(2\pi)^D}.
\label{eq:(3.5)}
\end{equation}
}
\begin{equation}
   K_t(x)_{\mu\nu}=\int_p\frac{\mathrm{e}^{ipx}}{p^2}
   \left[(\delta_{\mu\nu}p^2-p_\mu p_\nu)\mathrm{e}^{-tp^2}
   +p_\mu p_\nu \mathrm{e}^{-\alpha_0tp^2}\right]
\label{eq:(3.6)}
\end{equation}
is the heat kernel and
\begin{equation}
   R_\mu=2[B_\nu,\partial_\nu B_\mu]
   -[B_\nu,\partial_\mu B_\nu]
   +(\alpha_0-1)[B_\mu,\partial_\nu B_\nu]
   +[B_\nu,[B_\nu,B_\mu]]
\label{eq:(3.7)}
\end{equation}
denotes non-linear interaction terms. By iteratively solving
Eq.~\eqref{eq:(3.4)}, we have a perturbative expansion for the flowed
field~$B_\mu(t,x)$ in terms of the initial value~$A_\nu(y)$.

A similar flow may also be considered for fermion
fields~\cite{Luscher:2013cpa}. For our purpose, it is not necessary that the
flow of fermion fields is a gradient flow of the original fermion action. A
possible choice introduced in~Ref.~\cite{Luscher:2013cpa} is
\begin{align}
   &\partial_t\chi(t,x)=\left[\Delta-\alpha_0\partial_\mu B_\mu(t,x)\right]
   \chi(t,x),\qquad
   \chi(t=0,x)=\psi(x),
\label{eq:(3.8)}\\
   &\partial_t\Bar{\chi}(t,x)
   =\Bar{\chi}(t,x)
   \left[\overleftarrow{\Delta}
   +\alpha_0\partial_\mu B_\mu(t,x)\right],
   \qquad\Bar{\chi}(t=0,x)=\Bar{\psi}(x),
\label{eq:(3.9)}
\end{align}
where
\begin{align}
   &\Delta=D_\mu D_\mu,\qquad D_\mu=\partial_\mu+B_\mu,
\label{eq:(3.10)}\\
   &\overleftarrow{\Delta}=\overleftarrow{D}_\mu\overleftarrow{D}_\mu,
   \qquad\overleftarrow{D}_\mu\equiv\overleftarrow{\partial}_\mu-B_\mu.
\label{eq:(3.11)}
\end{align}
The formal solutions of the above flow equations are
\begin{align}
   \chi(t,x)&=\int\mathrm{d}^Dy\,
   \left[
   K_t(x-y)\psi(y)
   +\int_0^t\mathrm{d}s\,K_{t-s}(x-y)\Delta'\chi(s,y)
   \right],
\label{eq:(3.12)}\\
   \Bar{\chi}(t,x)&=\int\mathrm{d}^Dy\,
   \left[
   \Bar{\psi}(y)K_t(x-y)
   +\int_0^t\mathrm{d}s\,\Bar{\chi}(s,y)\overleftarrow{\Delta}'K_{t-s}(x-y)
   \right],
\label{eq:(3.13)}
\end{align}
where
\begin{equation}
   K_t(x)\equiv\int_p \mathrm{e}^{ipx}\mathrm{e}^{-tp^2}=\frac{\mathrm{e}^{-x^2/4t}}{(4\pi t)^{D/2}},
\label{eq:(3.14)}
\end{equation}
and
\begin{align}
   \Delta'&\equiv(1-\alpha_0)\partial_\mu B_\mu
   +2B_\mu\partial_\mu+B_\mu B_\mu,
\label{eq:(3.15)}\\
   \overleftarrow{\Delta}'&
   \equiv-(1-\alpha_0)\partial_\mu B_\mu
   -2\overleftarrow{\partial}_\mu B_\mu+B_\mu B_\mu.
\label{eq:(3.16)}
\end{align}

The initial values for the above flow, $A_\mu(x)$, $\psi(x)$,
and~$\Bar{\psi}(x)$, are quantum fields being subject to the functional
integral. The quantum correlation functions of the flowed fields are thus
obtained by expressing the flowed fields in terms of original un-flowed fields
(the initial values) and taking the functional average of the latter.
Equations~\eqref{eq:(3.4)}, \eqref{eq:(3.12)}, and~\eqref{eq:(3.13)} provide an
explicit method to carry this out. For example, in the lowest (tree-level)
approximation, we have
\begin{equation}
   \left\langle B_\mu^a(t,x)B_\nu^b(s,y)\right\rangle
   =g_0^2\delta^{ab}\delta_{\mu\nu}\int_p\mathrm{e}^{ip(x-y)}\frac{\mathrm{e}^{-(t+s)p^2}}{p^2},
\label{eq:(3.17)}
\end{equation}
in the ``Feynman gauge'' in which $\lambda_0=\alpha_0=1$, where $\lambda_0$ is
the conventional gauge-fixing parameter. Similarly, for the fermion field, in
the tree-level approximation,
\begin{equation}
   \left\langle\chi(t,x)\Bar{\chi}(s,y)\right\rangle
   =\int_p\mathrm{e}^{ip(x-y)}\frac{\mathrm{e}^{-(t+s)p^2}}{i\Slash{p}+m_0}.
\label{eq:(3.18)}
\end{equation}
% for the ghost field
% \begin{equation}
%    \left\langle c^a(x)\Bar{c}^b(y)\right\rangle_0
%    =g_0^2\delta^{ab}\int_p\mathrm{e}^{ip(x-y)}\frac{1}{p^2}.
% \end{equation}
Besides these ``quantum propagators,'' we also have heat kernels,
Eqs.~\eqref{eq:(3.6)} and~\eqref{eq:(3.14)}, in the perturbative expansion of
Eqs.~\eqref{eq:(3.4)}, \eqref{eq:(3.12)}, and~\eqref{eq:(3.13)}.

We now explain a diagrammatic representation of the perturbative expansion of
flowed fields (the flow Feynman diagram). For quantum
propagators~\eqref{eq:(3.17)} and~\eqref{eq:(3.18)}, we use the standard
convention that the free propagator of the gauge boson is denoted by a wavy
line and the free propagator of the fermion is denoted by an arrowed straight
line. We stick to these conventions because these are quite natural in a system
containing fermions. In~Refs.~\cite{Suzuki:2013gza} and~\cite{Luscher:2011bx},
on the other hand, the arrowed straight line was adopted to represent the gauge
boson heat kernel~\eqref{eq:(3.6)}. Since we already used this in this paper
for the fermion propagator, we instead use ``doubled lines'' to represent heat
kernels~\eqref{eq:(3.6)} and~\eqref{eq:(3.14)}. For instance,
Figs.~\ref{fig:1}--\ref{fig:5} are one-loop flow Feynman diagrams which
contribute to the two-point function of the flowed fermion field. In these
figures, the doubled straight line represents the fermion heat
kernel~\eqref{eq:(3.14)}; the arrow denotes the flow of the fermion number, not
the direction of the flow time. Similarly, in~Fig.~\ref{fig:15}, the doubled
wavy line is the gauge boson heat kernel~\eqref{eq:(3.6)}. In the present
representation, we thus lose the information of the direction of the flow time,
which is represented by an arrow in~Refs.~\cite{Suzuki:2013gza}
and~\cite{Luscher:2011bx}. This information, however, can readily be traced
back.

Another element of the flow Feynman diagram is the vertex. The vertices that
come from the original action~\eqref{eq:(2.1)} are denoted by filled circles,
while vertices that come through the flow equations, Eqs.~\eqref{eq:(3.7)},
\eqref{eq:(3.15)}, and~\eqref{eq:(3.16)}, are denoted by open circles as
in~Figs.~\ref{fig:1}--\ref{fig:5}; these conventions for the vertex are
identical to those of~Refs.~\cite{Suzuki:2013gza} and~\cite{Luscher:2011bx}.

\subsection{Ringed fermion fields}
A salient feature of the flowed fields is the UV finiteness: Any correlation
functions of the flowed gauge field $B_\mu(t,x)$ with strictly positive~$t$
are, when expressed in terms of renormalized parameters, UV finite
\emph{without\/} the multiplicative (wave function)
renormalization~\cite{Luscher:2011bx}. Moreover, this finiteness persists even
for the equal-point limit. Thus, any correlation functions of any local
products of~$B_\mu(t,x)$ (with $t>0$) are UV finite without further
renormalization other than the parameter renormalization. In other words,
although those local products are given by a certain combination of the bare
gauge field~$A_\mu(x)$ through the flow equation, they are renormalized finite
quantities. A basic reason for this UV finiteness is that the propagator of the
flowed gauge field~\eqref{eq:(3.17)} contains the Gaussian dumping
factor~$\sim \mathrm{e}^{-tp^2}$ which effectively provides a UV cutoff for~$t>0$. To
prove the above finiteness, however, one has to also utilize a
Becchi--Rouet--Stora (BRS) symmetry underlying the present system that is
\emph{inhomogeneous\/} with respect to the gauge
potential~\cite{Luscher:2011bx}.

Regrettably, the above finiteness in the first sense does not hold for the
flowed fermion field. It \emph{requires\/} the wave function renormalization.
Although its propagator~\eqref{eq:(3.18)} also possesses the dumping
factor~$\sim \mathrm{e}^{-tp^2}$, the BRS transformation is \emph{homogeneous\/} on the
fermion field (and on general matter fields) and the finiteness proof
in~Ref.~\cite{Luscher:2011bx} does not apply. In fact, computation of the
one-loop diagrams in~Figs.~\ref{fig:1}--\ref{fig:5} (and diagrams with opposite
arrows) shows that the wave function renormalization
\begin{equation}
   \chi_R(t,x)=Z_\chi^{1/2}\chi(t,x),\qquad
   \Bar{\chi}_R(t,x)=Z_\chi^{1/2}\Bar{\chi}(t,x),\qquad
   Z_\chi=1+\frac{g^2}{(4\pi)^2}C_2(R)3\frac{1}{\epsilon}+O(g^4)
\label{eq:(3.19)}
\end{equation}
makes correlation functions UV finite~\cite{Luscher:2013cpa}. Finiteness in the
above second sense still holds: Any correlation functions of any
local products of $\chi_R(t,x)$ and~$\Bar{\chi}_R(t,x)$ remain UV
finite~\cite{Luscher:2013cpa}.

\begin{figure}
\begin{minipage}{0.3\textwidth}
\begin{center}
\includegraphics[width=3cm,clip]{images/C01}
\caption{C01}
\label{fig:1}
\end{center}
\end{minipage}
\begin{minipage}{0.3\textwidth}
\begin{center}
\includegraphics[width=3cm,clip]{images/C02}
\caption{C02}
\label{fig:2}
\end{center}
\end{minipage}
\begin{minipage}{0.3\textwidth}
\begin{center}
\includegraphics[width=3cm,clip]{images/C03}
\caption{C03}
\label{fig:3}
\end{center}
\end{minipage}
\end{figure}

\begin{figure}
\begin{minipage}{0.3\textwidth}
\begin{center}
\includegraphics[width=3cm,clip]{images/C04}
\caption{C04}
\label{fig:4}
\end{center}
\end{minipage}
\begin{minipage}{0.3\textwidth}
\begin{center}
\includegraphics[width=3cm,clip]{images/C05}
\caption{C05}
\label{fig:5}
\end{center}
\end{minipage}
\end{figure}

Although the finiteness in the second sense is quite useful for our purpose, we
still need to incorporate the renormalization factor~$Z_\chi$
in~Eq.~\eqref{eq:(3.19)}. This introduces a complication to our problem,
because we have to find a matching factor between $Z_\chi$ in the dimensional
regularization and that in the lattice regularization.

One possible way to avoid this complication is to normalize the fermion fields
by the vacuum expectation value of the fermion kinetic operator:\footnote{In
the kinetic operator, the summation over $N_{\mathrm{f}}$ fermion flavors is understood.
In the first version of the present paper, we used the scalar condensation to
normalize the fermion fields. This choice causes another complication
associated with the massless fermion and the use of the kinetic operator seems
much more appropriate.}
\begin{align}
   \mathring{\chi}(t,x)
   &=\sqrt{\frac{-2\dim(R)N_{\mathrm{f}}}
   {(4\pi)^2t^2
   \left\langle\Bar{\chi}(t,x)\overleftrightarrow{\Slash{D}}\chi(t,x)
   \right\rangle}}
   \,\chi(t,x),
\label{eq:(3.20)}\\
   \mathring{\Bar{\chi}}(t,x)
   &=\sqrt{\frac{-2\dim(R)N_{\mathrm{f}}}
   {(4\pi)^2t^2
   \left\langle\Bar{\chi}(t,x)\overleftrightarrow{\Slash{D}}\chi(t,x)
   \right\rangle}}
   \,\Bar{\chi}(t,x),
\label{eq:(3.21)}
\end{align}
where
\begin{equation}
   \overleftrightarrow{D}_\mu\equiv D_\mu-\overleftarrow{D}_\mu,
\label{eq:(3.22)}
\end{equation}
so that the multiplicative renormalization factor~$Z_\chi$ is cancelled out in
the new ringed variables. Note that the mass dimension
of~$\mathring{\chi}(t,x)$ and~$\mathring{\Bar{\chi}}(t,x)$ is $3/2$ for
any~$D$, while that of~$\chi(t,x)$ and~$\Bar{\chi}(t,x)$ is~$(D-1)/2$.

The vacuum expectation value of the kinetic operator in the lowest (one-loop)
order approximation is given by diagram~D01 in~Fig.~\ref{fig:6}. For
$D=4-2\epsilon$ dimensions, we have
\begin{equation}
   \left\langle\Bar{\chi}(t,x)\overleftrightarrow{\Slash{D}}\chi(t,x)
   \right\rangle
   =\frac{-2\dim(R)N_{\mathrm{f}}}{(4\pi)^2t^2}(8\pi t)^\epsilon\left[1+O(m_0^2t)\right].
\label{eq:(3.23)}
\end{equation}
Note that the mass scale, which is required for the vacuum expectation value,
is supplied by the flow time~$t$ in the present setup. Having obtained this
expression, the constant factors in~Eqs.~\eqref{eq:(3.20)}
and~\eqref{eq:(3.21)} have been chosen so that the difference between the
ringed and the original variables becomes~$O(g_0^2)$ for sufficiently small
flow time ($m_0^2t\ll1$).

\begin{figure}
\begin{center}
\includegraphics[width=3cm,clip]{images/D01}
\caption{D01}
\label{fig:6}
\end{center}
\end{figure}

\begin{figure}
\begin{minipage}{0.3\textwidth}
\begin{center}
\includegraphics[width=3cm,clip]{images/D02}
\caption{D02}
\label{fig:7}
\end{center}
\end{minipage}
\begin{minipage}{0.3\textwidth}
\begin{center}
\includegraphics[width=3cm,clip]{images/D03}
\caption{D03}
\label{fig:8}
\end{center}
\end{minipage}
\begin{minipage}{0.3\textwidth}
\begin{center}
\includegraphics[width=3cm,clip]{images/D04}
\caption{D04}
\label{fig:9}
\end{center}
\end{minipage}
\end{figure}

\begin{figure}
\begin{minipage}{0.3\textwidth}
\begin{center}
\includegraphics[width=3cm,clip]{images/D05}
\caption{D05}
\label{fig:10}
\end{center}
\end{minipage}
\begin{minipage}{0.3\textwidth}
\begin{center}
\includegraphics[width=3cm,clip]{images/D06}
\caption{D06}
\label{fig:11}
\end{center}
\end{minipage}
\begin{minipage}{0.3\textwidth}
\begin{center}
\includegraphics[width=3cm,clip]{images/D07}
\caption{D07}
\label{fig:12}
\end{center}
\end{minipage}
\end{figure}

\begin{figure}
\begin{center}
\includegraphics[width=3cm,clip]{images/D08}
\caption{D08}
\label{fig:13}
\end{center}
\end{figure}

\begin{table}
\caption{Contribution of each diagram to Eq.~\eqref{eq:(3.24)} in units
of~$\frac{-2\dim(R)N_{\mathrm{f}}}{(4\pi)^2t^2}\frac{g_0^2}{(4\pi)^2}C_2(R)$.}
\label{table:1}
\begin{center}
\renewcommand{\arraystretch}{2.2}
\setlength{\tabcolsep}{20pt}
\begin{tabular}{cr}
\toprule
 diagram & \\
\midrule
 D02 & $-\dfrac{1}{\epsilon}-2\ln(8\pi t)+O(m_0^2t)$ \\
 D03 & $2\dfrac{1}{\epsilon}+4\ln(8\pi t)+2+4\ln2-2\ln3+O(m_0^2t)$ \\
 D04 & $-20\ln2+16\ln3+O(m_0^2t)$ \\
 D05 & $12\ln2-5\ln3+O(m_0^2t)$ \\
 D06 & $-4\dfrac{1}{\epsilon}-8\ln(8\pi t)-2+O(m_0^2t)$ \\
 D07 & $8\ln2-4\ln3+O(m_0^2t)$ \\
 D08 & $-2\ln3+O(m_0^2t)$ \\
\bottomrule
\end{tabular}
\end{center}
\end{table}

The next-to-leading-order (i.e., two-loop) expression for the expectation
value~\eqref{eq:(3.23)} is given by the flow Feynman diagrams
in~Figs.~\ref{fig:7}--\ref{fig:13} (and diagrams with arrows with the opposite
direction). The computation of these diagrams is somewhat complicated but can
be completed in a similar manner to the calculation in~Appendix~B
of~Ref.~\cite{Luscher:2010iy} [with the integration formulas in our
Appendix~\ref{sec:B}, Eqs.~\eqref{eq:(B1)} and~\eqref{eq:(B2)}]. The
contribution of each diagram is tabulated in~Table~\ref{table:1}. In total, we
have
\begin{align}
   &\left\langle\Bar{\chi}(t,x)\overleftrightarrow{\Slash{D}}\chi(t,x)
   \right\rangle
\notag\\
   &=\frac{-2\dim(R)N_{\mathrm{f}}}{(4\pi)^2t^2}
   \left\{(8\pi t)^\epsilon
   +\frac{g_0^2}{(4\pi)^2}C_2(R)
   \left[-3\frac{1}{\epsilon}-6\ln(8\pi t)+\ln(432)\right]
   +O(m_0^2t)+O(g_0^4)\right\}.
\label{eq:(3.24)}
\end{align}
Using Eq.~\eqref{eq:(A1)} for the normalization factor
in~Eqs.~\eqref{eq:(3.20)} and~\eqref{eq:(3.21)}, we have
\begin{equation}
   \frac{-2\dim(R)N_{\mathrm{f}}}
   {(4\pi)^2t^2\left\langle\Bar{\chi}(t,x)
   \overleftrightarrow{\Slash{D}}\chi(t,x)\right\rangle}
   =Z(\epsilon)
   \left\{1+\frac{g^2}{(4\pi)^2}C_2(R)\left[3\frac{1}{\epsilon}-\Phi(t)\right]
   +O(m^2t)+O(g^4)\right\},
\label{eq:(3.25)}
\end{equation}
where
\begin{equation}
   Z(\epsilon)\equiv\frac{1}{(8\pi t)^\epsilon},
\label{eq:(3.26)}
\end{equation}
and
\begin{equation}
   \Phi(t)\equiv-3\ln(8\pi\mu^2t)+\ln(432).
\label{eq:(3.27)}
\end{equation}
